% GENERATED FILE. Edit canonical lessons, then run npm run generate:books.
% canonical-source-hash: fnv1a64:5613f57fd6814157
% canonical-lessons: AR-C190-mughanni, AR-C190-laib, AR-C190-mudarrib, AR-C190-riyadiyyat, AR-C190-jughrafiya

\chapter{Singer, Player, Coach, Mathematics, Geography}
\label{ch:ar-singer-player-coach-mathematics-geography}
\hlchaptermodality{190}

\begin{chapteropening}
\textbf{By the end of this chapter:} \emph{I can say a singer, a player, a coach, mathematics and geography.}

Complete the last lesson of chapter 190: I can say a singer, a player, a coach, mathematics and geography.
\end{chapteropening}

\section[\texorpdfstring{\ar{مغني} (mughannī)}{مغني (mughannī)}]{\ar{مغني} (mughannī) — a singer}

\label{lesson:AR-C190-mughanni}



\begin{quote}
Before the new one: say the Arabic for a translator, then the Arabic for an interview.
\end{quote}

\subsection*{You'll want to know: \ar{مغني}}
\textbf{\ar{مغني}} — \emph{mughannī} — \textquotedblleft{}a singer\textquotedblright{}.

\begin{cousinweb}
\textbf{\ar{مغني}} (\emph{mughannī}) is built on the native Arabic root \emph{gh-n-y}, \textquotedblleft{}to be free from want\textquotedblright{} (Lane), whose Form II \emph{ghannā} means \textquotedblleft{}to sing\textquotedblright{}. \emph{Mughannī} is the Form II participle, \textquotedblleft{}one who sings\textquotedblright{}.
\end{cousinweb}

\begin{grammarlens}[title={What you've built}]
One more word for this chapter.
\end{grammarlens}

\subsection*{Your turn}
\begin{itemize}
\raggedright
  \item \emph{Say it:} \emph{mughannī}
  \item \emph{Say it:} \emph{mughannī}, once more
  \item \emph{Say it:} say \emph{mughannī}
  \item \emph{Recall:} say the Arabic for a translator, then the Arabic for an interview, then say \emph{mughannī} again
\end{itemize}

\begin{tcolorbox}[breakable,colback=teal!4,colframe=teal!35!black,title={Before you move on}]
What does \ar{مغني} mean? (\textquotedblleft{}A singer\textquotedblright{}.) Say it once more.
\end{tcolorbox}
\section[\texorpdfstring{\ar{لاعب} (lāʿib)}{لاعب (lāʿib)}]{\ar{لاعب} (lāʿib) — a player}

\label{lesson:AR-C190-laib}



\begin{quote}
Before the new one: say the Arabic for an interview, then the Arabic for a singer.
\end{quote}

\subsection*{You'll want to know: \ar{لاعب}}
\textbf{\ar{لاعب}} — \emph{lāʿib} — \textquotedblleft{}a player\textquotedblright{}.

\begin{cousinweb}
\textbf{\ar{لاعب}} (\emph{lāʿib}) is built on the native Arabic root \emph{l-ʿ-b}, \textquotedblleft{}to play\textquotedblright{} (Lane). \emph{Lāʿib} is the active participle, \textquotedblleft{}one who plays\textquotedblright{}.
\end{cousinweb}

\begin{grammarlens}[title={What you've built}]
One more word for this chapter.
\end{grammarlens}

\subsection*{Your turn}
\begin{itemize}
\raggedright
  \item \emph{Say it:} \emph{lāʿib}
  \item \emph{Say it:} \emph{lāʿib}, once more
  \item \emph{Say it:} say \emph{lāʿib}
  \item \emph{Recall:} say the Arabic for an interview, then the Arabic for a singer, then say \emph{lāʿib} again
\end{itemize}

\begin{tcolorbox}[breakable,colback=teal!4,colframe=teal!35!black,title={Before you move on}]
What does \ar{لاعب} mean? (\textquotedblleft{}A player\textquotedblright{}.) Say it once more.
\end{tcolorbox}
\section[\texorpdfstring{\ar{مدرب} (mudarrib)}{مدرب (mudarrib)}]{\ar{مدرب} (mudarrib) — a coach, a trainer}

\label{lesson:AR-C190-mudarrib}



\begin{quote}
Before the new one: say the Arabic for a singer, then the Arabic for a player.
\end{quote}

\subsection*{You'll want to know: \ar{مدرب}}
\textbf{\ar{مدرب}} — \emph{mudarrib} — \textquotedblleft{}a coach, a trainer\textquotedblright{}.

\begin{cousinweb}
\textbf{\ar{مدرب}} (\emph{mudarrib}) is built on the native Arabic root \emph{d-r-b}, \textquotedblleft{}to be trained, to be accustomed\textquotedblright{} (Lane). \emph{Mudarrib} is the Form II participle, \textquotedblleft{}one who trains\textquotedblright{}.
\end{cousinweb}

\begin{grammarlens}[title={What you've built}]
One more word for this chapter.
\end{grammarlens}

\subsection*{Your turn}
\begin{itemize}
\raggedright
  \item \emph{Say it:} \emph{mudarrib}
  \item \emph{Say it:} \emph{mudarrib}, once more
  \item \emph{Say it:} say \emph{mudarrib}
  \item \emph{Recall:} say the Arabic for a singer, then the Arabic for a player, then say \emph{mudarrib} again
\end{itemize}

\begin{tcolorbox}[breakable,colback=teal!4,colframe=teal!35!black,title={Before you move on}]
What does \ar{مدرب} mean? (\textquotedblleft{}A coach, a trainer\textquotedblright{}.) Say it once more.
\end{tcolorbox}
\section[\texorpdfstring{\ar{رياضيات} (riyāḍiyyāt)}{رياضيات (riyāḍiyyāt)}]{\ar{رياضيات} (riyāḍiyyāt) — mathematics}

\label{lesson:AR-C190-riyadiyyat}



\begin{quote}
Before the new one: say the Arabic for a player, then the Arabic for a coach.
\end{quote}

\subsection*{You'll want to know: \ar{رياضيات}}
\textbf{\ar{رياضيات}} — \emph{riyāḍiyyāt} — \textquotedblleft{}mathematics\textquotedblright{}.

\begin{cousinweb}
\textbf{\ar{رياضيات}} (\emph{riyāḍiyyāt}) is built on the native Arabic root \emph{r-w-ḍ}, \textquotedblleft{}to train, to tame\textquotedblright{} (Lane). \emph{Riyāḍiyyāt} is built on \emph{riyāḍa}, \textquotedblleft{}training, exercise\textquotedblright{}, with the plural abstract ending \emph{-iyyāt}.
\end{cousinweb}

\begin{grammarlens}[title={What you've built}]
One more word for this chapter.
\end{grammarlens}

\subsection*{Your turn}
\begin{itemize}
\raggedright
  \item \emph{Say it:} \emph{riyāḍiyyāt}
  \item \emph{Say it:} \emph{riyāḍiyyāt}, once more
  \item \emph{Say it:} say \emph{riyāḍiyyāt}
  \item \emph{Recall:} say the Arabic for a player, then the Arabic for a coach, then say \emph{riyāḍiyyāt} again
\end{itemize}

\begin{tcolorbox}[breakable,colback=teal!4,colframe=teal!35!black,title={Before you move on}]
What does \ar{رياضيات} mean? (\textquotedblleft{}Mathematics\textquotedblright{}.) Say it once more.
\end{tcolorbox}
\section[\texorpdfstring{\ar{جغرافيا} (jughrāfiyā)}{جغرافيا (jughrāfiyā)}]{\ar{جغرافيا} (jughrāfiyā) — geography}

\label{lesson:AR-C190-jughrafiya}



\begin{quote}
Before the new one: say the Arabic for a coach, then the Arabic for mathematics.
\end{quote}

\subsection*{You'll want to know: \ar{جغرافيا}}
\textbf{\ar{جغرافيا}} — \emph{jughrāfiyā} — \textquotedblleft{}geography\textquotedblright{}.

\begin{cousinweb}
\textbf{\ar{جغرافيا}} (\emph{jughrāfiyā}) is a loan from Greek \emph{geōgraphia}, \textquotedblleft{}description of the earth\textquotedblright{} (EtymDB); it has no Arabic root.
\end{cousinweb}

\begin{grammarlens}[title={What you've built}]
One more word for this chapter.
\end{grammarlens}

\subsection*{Your turn}
\begin{itemize}
\raggedright
  \item \emph{Say it:} \emph{jughrāfiyā}
  \item \emph{Say it:} \emph{jughrāfiyā}, once more
  \item \emph{Say it:} say \emph{jughrāfiyā}
  \item \emph{Recall:} say the Arabic for a coach, then the Arabic for mathematics, then say \emph{jughrāfiyā} again
\end{itemize}

\begin{tcolorbox}[breakable,colback=teal!4,colframe=teal!35!black,title={Before you move on}]
What does \ar{جغرافيا} mean? (\textquotedblleft{}Geography\textquotedblright{}.) Say it once more.
\end{tcolorbox}
